\chapter{器官在工作}\label{ch:g7:organs-at-work}

爬上三段楼梯，清点身体的反应：双腿
像在燃烧，\omterm{def:g4:heart-and-blood:heart}{心脏}猛烈搏动，\omterm{def:g7:what-is-respiration:respiration}{呼吸}急促——奇怪的是，
\omterm{def:g1:the-five-senses:sense}{皮肤}也发红发热。没有人下令让这些变化发生，
它们是自动发生的，而且比例恰到好处。本章测量
工作中的器官究竟消耗多少，
以及工作一开始，身体便如何重新调配供给。

\section{工作肌肉的消耗}

\begin{proposition}[肌肉的消耗账单]\label{prop:g7:organs-at-work:bill}
工作中的\omterm{def:g3:bones-and-muscles:muscle}{肌肉}在消耗，而这种消耗是可以测量的。
与同一块\omterm{def:g3:bones-and-muscles:muscle}{肌肉}静止时相比，剧烈工作的
\omterm{def:g3:bones-and-muscles:muscle}{肌肉}：
\begin{itemize}
  \item 从流经\omterm{def:g3:bones-and-muscles:muscle}{肌肉}的\omterm{prop:g4:journey-of-food:absorb}{血液}中摄取多得多的\emph{\omterm{def:g7:what-is-respiration:respiration}{氧气}}
        ——为静止时摄取量的数倍；
  \item 摄取多得多的\emph{葡萄糖}\index{葡萄糖}——
        即\omterm{prop:g4:journey-of-food:absorb}{血液}中现成的糖类养料，是其日常燃料；
  \item 相应地释放更多的\emph{\omterm{def:g7:what-is-respiration:respiration}{二氧化碳}}
        ——还有\emph{热量}。
\end{itemize}
这正是\cref{prop:g7:what-is-respiration:power}所说的\omterm{def:g5:first-look-at-cells:cell}{细胞}燃烧，
从器官的“仪表”上看：燃料与\omterm{def:g7:what-is-respiration:respiration}{氧气}输入，
做功，同时排出废气和热量。
\end{proposition}
\admitted

\begin{example}[肌肉两端的测量]\label{ex:g7:organs-at-work:measure}
生理学家比较\emph{流入}\omterm{def:g3:bones-and-muscles:muscle}{肌肉}的\omterm{prop:g4:journey-of-food:absorb}{血液}与
\emph{流出}\omterm{def:g3:bones-and-muscles:muscle}{肌肉}的\omterm{prop:g4:journey-of-food:absorb}{血液}。静止时，流出的\omterm{prop:g4:journey-of-food:absorb}{血液}
只释放了一小部分\omterm{def:g7:what-is-respiration:respiration}{氧气}和\omterm{prop:g7:organs-at-work:bill}{葡萄糖}。剧烈工作时，
同样的比较显示流出的\omterm{prop:g4:journey-of-food:absorb}{血液}被大幅榨取，
——大部分可用\omterm{def:g7:what-is-respiration:respiration}{氧气}已被取走，而\omterm{def:g7:what-is-respiration:respiration}{二氧化碳}
则大幅增加。同一器官、同样的仪表，变化的只是工作，
而消耗的账单也随之改变。
\end{example}

\begin{remark}[热量即凭证]\label{rem:g7:organs-at-work:warmth}
工作中的身体所产生的热量并非偶然：燃料燃烧
总会释放热量，而\omterm{def:g3:bones-and-muscles:muscle}{肌肉}就是身体的“锅炉房”。
寒战便是刻意而为的把戏——让\omterm{def:g3:bones-and-muscles:muscle}{肌肉}
仅为产热而运动，
（见\cref{ex:g3:balanced-diet:rest}中的静息能耗，短暂
翻倍）。冲刺时\omterm{def:g1:the-five-senses:sense}{皮肤}发红，是散热系统
在把“锅炉房”的热量排出体外，这将在下一节讨论。
\end{remark}

\section{按需供给}

\begin{proposition}[身体重新调配供给]\label{prop:g7:organs-at-work:supply}
工作一开始，供给线路便重新调整，以支付
这笔“账单”（\cref{ex:g4:heart-and-blood:effort}已预告过；
此处给出完整的调配方案）：
\begin{enumerate}
  \item \emph{\omterm{def:g7:what-is-respiration:respiration}{呼吸}}加深、加快——
        使\omterm{prop:g4:journey-of-food:absorb}{血液}在肺部更快地重新装载\omterm{def:g7:what-is-respiration:respiration}{氧气}
        （见\cref{ex:g4:breathing:running}）；
  \item \emph{\omterm{def:g4:heart-and-blood:heart}{心脏}}搏动更快更有力——每分钟
        搏动的次数增多
        （见\cref{prop:g4:heart-and-blood:pulse}）；
  \item \omterm{prop:g4:journey-of-food:absorb}{血液}被\emph{重新分配}：为工作\omterm{def:g3:bones-and-muscles:muscle}{肌肉}
        （及散热的\omterm{def:g1:the-five-senses:sense}{皮肤}）供血的\omterm{def:g4:heart-and-blood:vessels}{血管}扩张，
        而静息中的\omterm{def:g4:journey-of-food:tract}{消化道}\omterm{def:g4:heart-and-blood:vessels}{血管}收缩——
        同一批\omterm{prop:g4:journey-of-food:absorb}{血液}被改道流向
        正在耗能的部位。
\end{enumerate}
三者都随工作自动发生，由
\omterm{def:g5:brain-in-command:system}{大脑}在无须指令的情况下调控
（见\cref{rem:g5:brain-in-command:more}）。
\end{proposition}
\admitted

\begin{example}[从内部感受到的调配]\label{ex:g7:organs-at-work:felt}
逐条解释爬楼梯冲刺时的身体“清点”：
\omterm{def:g7:what-is-respiration:respiration}{呼吸}急促——是在补充\omterm{def:g7:what-is-respiration:respiration}{氧气}；心跳如擂鼓——
搏动加快；双腿灼痛——是\omterm{def:g3:bones-and-muscles:muscle}{肌肉}的“账单”正以
极限速率支付；\omterm{def:g1:the-five-senses:sense}{皮肤}发红发热——是把热量
送到散热的体表。连经典的“岔气”
与饱餐后的困倦也都能解释：\omterm{def:g4:journey-of-food:digestion}{消化}与冲刺
争夺改道后的\omterm{prop:g4:journey-of-food:absorb}{血液}，这正是
\cref{met:g4:journey-of-food:habits}曾提醒
不要在饱胃时剧烈运动。
\end{example}

\begin{omfigure}
\begin{tikzpicture}
\begin{axis}[omaxis,
    width=9.6cm, height=5.8cm,
    xlabel={分钟}, ylabel={每分钟心跳次数},
    xmin=0, xmax=13.6, ymin=40, ymax=200,
    xtick={0,2,4,6,8,10,12},
    ytick={60,100,140,180},
    ]
  \addplot[very thick, omThm, smooth] coordinates {
    (0,72) (1,74) (2,120) (3,150) (4,158) (5,160)
    (6,120) (7,100) (8,88) (9,80) (10,76) (11,74) (12,73)
  };
  \draw[dashed, omExo] (axis cs:2,40) -- (axis cs:2,180);
  \draw[dashed, omExo] (axis cs:5,40) -- (axis cs:5,180);
  \node[font=\footnotesize] at (axis cs:1.0,52) {静息};
  \node[font=\footnotesize] at (axis cs:3.5,52) {用力};
  \node[font=\footnotesize] at (axis cs:8.6,52) {恢复};
\end{axis}
\end{tikzpicture}

{\small 三分钟剧烈用力中的\omterm{prop:g4:heart-and-blood:pulse}{脉搏}变化：攀登
时随工作上升，接近最高值时趋于平缓，
以及缓慢回落——即\omterm{def:g7:organs-at-work:recovery}{恢复}阶段——身体
重新结算其“账目”。}
\end{omfigure}

\begin{definition}[恢复]\label{def:g7:organs-at-work:recovery}
用力停止时，\omterm{prop:g4:heart-and-blood:pulse}{脉搏}和\omterm{def:g7:what-is-respiration:respiration}{呼吸}并不随之停止：它们
只是逐渐回落——即
\emph{恢复}\index{恢复}阶段。身体正在偿还
欠账：补充\omterm{def:g7:what-is-respiration:respiration}{氧气}储备，清除用力产生的
废物，散发积存的热量。体能越好的身体，
恢复越快，这也正是为什么恢复
时间本身就是一项体能测量指标。
\end{definition}

\section{训练与机器的极限}

\begin{proposition}[训练提高上限]\label{prop:g7:organs-at-work:training}
经常锻炼，整条供应链都会增强
（见\cref{rem:g4:heart-and-blood:rest}，
\cref{met:g4:breathing:care}）：\omterm{def:g4:heart-and-blood:heart}{心脏}每次搏动输送的\omterm{prop:g4:journey-of-food:absorb}{血液}更多，
\omterm{def:g7:what-is-respiration:respiration}{呼吸}肌力量增强，\omterm{def:g3:bones-and-muscles:muscle}{肌肉}
自身变大并储存更多燃料，微\omterm{def:g4:heart-and-blood:vessels}{血管}
数量增多。结果是：同样的用力，对受过训练的
身体来说\omterm{prop:g4:heart-and-blood:pulse}{脉搏}更低，其上限——即它能承担的
最强用力——随之提高。机器通过
自身的使用得以重建。
\end{proposition}
\admitted

\begin{example}[两名跑者，同一座山]\label{ex:g7:organs-at-work:runners}
未经训练者与受过训练者跑同一座山。未经训练者：\omterm{prop:g4:heart-and-blood:pulse}{脉搏}
中途达 $180$，气喘吁吁，被迫步行——供应
线路承担不起这笔开销，欠账迫使其停下。
受过训练者：山顶时\omterm{prop:g4:heart-and-blood:pulse}{脉搏} $150$，能够说话，\omterm{def:g7:organs-at-work:recovery}{恢复}
只需三分钟。同一座山，同一笔开销——
上限不同，是数月来像这样循序渐进地
攀登这类山坡所建立的。
\end{example}

\begin{method}[读懂自己的机器]\label{met:g7:organs-at-work:read}
用手表进行的自我实验
（测\omterm{prop:g4:heart-and-blood:pulse}{脉搏}的方法见\cref{met:g4:heart-and-blood:pulse}
）：
\begin{enumerate}
  \item 坐着、平静时，记录静息\omterm{prop:g4:heart-and-blood:pulse}{脉搏}；
  \item 持续用力三分钟——上下踏台阶、
        跳绳；
  \item 立即记录\omterm{prop:g4:heart-and-blood:pulse}{脉搏}，此后每分钟记录一次，直到
        \omterm{def:g7:organs-at-work:recovery}{恢复}静息：即\omterm{def:g7:organs-at-work:recovery}{恢复}时间；
  \item 在一学期的规律锻炼中每周重复一次：
        观察静息\omterm{prop:g4:heart-and-blood:pulse}{脉搏}、用力时\omterm{prop:g4:heart-and-blood:pulse}{脉搏}与\omterm{def:g7:organs-at-work:recovery}{恢复}
        时间，都逐渐下降——这正是
        \cref{prop:g7:organs-at-work:training}所述的重建在你的
        自己的数据中的体现。
\end{enumerate}
\end{method}

\section{练习}

\begin{exercise}[$\star$]\label{exo:g7:organs-at-work:1}
列出工作时的\omterm{def:g3:bones-and-muscles:muscle}{肌肉}摄取更多的哪些物质，
以及释放更多的哪些物质。
\end{exercise}

\begin{exercise}[$\star$]\label{exo:g7:organs-at-work:2}
什么是\omterm{prop:g7:organs-at-work:bill}{葡萄糖}，工作的\omterm{def:g3:bones-and-muscles:muscle}{肌肉}从何处获得它？
\end{exercise}

\begin{exercise}[$\star$]\label{exo:g7:organs-at-work:3}
描述流入与流出的\omterm{prop:g4:journey-of-food:absorb}{血液}之比较，以及
它在静息与用力时分别说明了什么。
\end{exercise}

\begin{exercise}[$\star$]\label{exo:g7:organs-at-work:4}
说出三种供应重组方式，
（见\cref{prop:g7:organs-at-work:supply}）。
\end{exercise}

\begin{exercise}[$\star$]\label{exo:g7:organs-at-work:5}
为什么剧烈劳动会使你发热，以及用本章的术语
说明什么是寒战？
\end{exercise}

\begin{exercise}[$\star$]\label{exo:g7:organs-at-work:6}
什么是\omterm{def:g7:organs-at-work:recovery}{恢复}？\omterm{def:g7:organs-at-work:recovery}{恢复}期间身体正在偿还哪些
欠账？
\end{exercise}

\begin{exercise}[$\star\star$]\label{exo:g7:organs-at-work:7}
读\omterm{prop:g4:heart-and-blood:pulse}{脉搏}图：静息值、平台期数值，以及
大约需要多长时间\omterm{def:g7:organs-at-work:recovery}{恢复}。训练会如何改变
其中每一项？
\end{exercise}

\begin{exercise}[$\star\star$]\label{exo:g7:organs-at-work:8}
用血流重新分配规则解释午饭后的乏力和短跑者的
岔气。
\end{exercise}

\begin{exercise}[$\star\star$]\label{exo:g7:organs-at-work:9}
为什么用力时\omterm{def:g1:the-five-senses:sense}{皮肤}潮红是一种\emph{输送}
改变而非故障？它解决的是谁的问题？
\end{exercise}

\begin{exercise}[$\star\star$]\label{exo:g7:organs-at-work:10}
列出训练提高上限所依赖的三种重建，
以及每一种的日常表现。
\end{exercise}

\begin{exercise}[$\star\star$]\label{exo:g7:organs-at-work:11}
两名学生按\cref{met:g7:organs-at-work:read}
进行测量：静息\omterm{prop:g4:heart-and-blood:pulse}{脉搏}为 $68$ 和 $86$；\omterm{def:g7:organs-at-work:recovery}{恢复}时间
分别为三分钟和八分钟。哪一台是受过训练的机器，以及
这两个数字各自确切证明了什么？
\end{exercise}

\begin{exercise}[$\star\star\star$]\label{exo:g7:organs-at-work:12}
用一段话把三章串联起来：\omterm{def:g3:bones-and-muscles:muscle}{肌肉}的“账单”
（\cref{prop:g7:organs-at-work:bill}）、\omterm{def:g7:what-is-respiration:respiration}{呼吸}的目的
（\cref{prop:g7:what-is-respiration:power}），以及\omterm{prop:g4:journey-of-food:absorb}{血液}的
输送服务
（\cref{prop:g4:heart-and-blood:carries}）——合起来，作为对一次短跑的
完整叙述。
\end{exercise}

\section{问题：学校体能研究}

\begin{problem}[{周末问题——一学期的脉搏数据，像生理学家那样解读，
读懂这些数据}]\label{pb:g7:organs-at-work:1}
全班在一学期内每周一都按\cref{met:g7:organs-at-work:read}测量一次：
静息\omterm{prop:g4:heart-and-blood:pulse}{脉搏}、做三分钟台阶练习后的\omterm{prop:g4:heart-and-blood:pulse}{脉搏}，
以及\omterm{def:g7:organs-at-work:recovery}{恢复}时间。两份匿名记录：
学生 A——第 1 周：静息 $84$，运动时 $168$，
\omterm{def:g7:organs-at-work:recovery}{恢复} $9$ 分钟；第 12 周：静息 $72$，运动时 $150$，
\omterm{def:g7:organs-at-work:recovery}{恢复} $5$ 分钟。学生 B——第 1 周：静息 $70$，
运动时 $150$，\omterm{def:g7:organs-at-work:recovery}{恢复} $4$ 分钟；第 12 周：静息 $69$，
运动时 $148$，\omterm{def:g7:organs-at-work:recovery}{恢复} $4$ 分钟。

\medskip\noindent\textbf{第一部分——单次课中的情况。}
\begin{enumerate}
  \item 在台阶练习期间，列出大腿\omterm{def:g3:bones-and-muscles:muscle}{肌肉}的
        \omterm{def:g7:what-is-respiration:respiration}{氧气}与\omterm{prop:g7:organs-at-work:bill}{葡萄糖}摄取量的变化，以及其
        \omterm{def:g7:what-is-respiration:respiration}{二氧化碳}与热量产生的变化。
  \item 解释为什么在用力的几分钟内，\omterm{def:g7:what-is-respiration:respiration}{呼吸}与\omterm{prop:g4:heart-and-blood:pulse}{脉搏}
        \emph{同时}上升。
  \item 一名学生在第三分钟时身上发热、面色潮红。
        请给出“产热炉与散热”的解释。
  \item 为什么台阶练习刚一停止，\omterm{prop:g4:heart-and-blood:pulse}{脉搏}
        并不会立刻回落到静息水平？
\end{enumerate}

\medskip\noindent\textbf{第二部分——解读整个学期的数据。}
\begin{enumerate}[resume]
  \item 描述学生 A 整个学期的变化，并说出
        每个数值背后依赖的
        \cref{prop:g7:organs-at-work:training}中的
        哪些重建。
  \item 学生 B 的数字几乎没有变化。给出两种
        合理的解读，以及学生 B 课外生活中的哪一个
        额外事实可以判定二者中哪一种成立。
  \item 第 5 周，一名学生在午饭匆忙吃完后测量，
        得到了古怪的数值。在
        \cref{prop:g7:organs-at-work:supply}中，哪一条规则受到干扰，
        以及测量方案应当就
        测量时间规定什么？
  \item 为什么每名学生都必须始终以相同的方式测量——
        即用力强度相同、计时方式相同？
        （方法中的一位“老朋友”，请说出它的名称。）
\end{enumerate}

\medskip\noindent\textbf{第三部分——生理学家的
来信。}
\begin{enumerate}[resume]
  \item 为学生 A 写一段两句话的学期小结，
        供班级公告栏使用，须用到\emph{上限}和
        \emph{\omterm{def:g7:organs-at-work:recovery}{恢复}}。
  \item 一位家长问静息\omterm{prop:g4:heart-and-blood:pulse}{脉搏}偏低是否令人担忧。
        请借用
        \cref{exo:g4:heart-and-blood:11}中的骑车人作答。
  \item 全班希望下学期跟踪一个指标。
        请论证应以\omterm{def:g7:organs-at-work:recovery}{恢复}时间而非运动时\omterm{prop:g4:heart-and-blood:pulse}{脉搏}为指标，
        并运用\cref{def:g7:organs-at-work:recovery}。
  \item 用一句话总结本研究的寓意：
        是什么重建了学生 A 的“机器”？
\end{enumerate}
\end{problem}
